% Français 9115, batch 10 — Gallica views 181–200.
% Pass: Opus 5 (claude-opus-5), 2026-09-15 — first pass, unchecked against
% the views by a human.
%
% What the twenty views hold. One thing, without a break: the fair copy of
% Euler's Latin memoir « Investigatio motuum quibus laminæ et virgæ elasticæ
% contremiscunt » (Acta Petropolitana pro anno 1779, Pars prior, as its first
% leaf says at view 144 in batch 8) runs on through the whole batch, in the
% same regular hand as batch 8, with the copy's parenthesised page numbers
% (140) at view 181 to (156) at view 200. Batch 9, written in parallel,
% ends at view 180 on page (139), the Scholion § 34 of case III, with the
% ink number 89; view 181 opens on the title of case IV and page (140), so
% the seam is clean. Seventeen views carry text.
% Three are not transcribed because they photograph a page a second time:
% view 191 is page (149) again (view 190), and views 196 and 197 are pages
% (152) and (153) again (views 194 and 195) — same text, same blots, same
% ink numbers, the leaf only framed differently. There are no blank views.
%
%   v. 181–183  « Evolutio casus IV », both ends simply fixed by a pin
%               (Problema § 35): $\beta = -\alpha$, $\delta = 0$, $\alpha = 0$,
%               $\sin\omega = 0$, so $\omega = \pi, 2\pi, 3\pi \dots$ and the
%               sounds rise as the squares 1, 4, 9, 16, named as notes $C$,
%               $\bar c$, $\bar{\bar d}$ …; the general motion as a sum of
%               sines of $\pi u$, $2\pi u$ …. Corollarium 1, § 36: the
%               fundamental takes the figure of a simply vibrating string
%               (« Tab. III fig. 1 », « fig. 2 »), one more node for each
%               higher sound. § 37: the four fundamentals compared, gis, d,
%               C, fis, with $\frac94$, $\frac{25}{16}$, $\frac{9}{25}$, 1.
%               Scholion § 38: the only case where every $\omega$ is exact,
%               hence sounds purer than a string's.
%   v. 184–187  « Evolutio casus V », one end pinned, the other built into a
%               wall (Problema § 39): the same equation
%               $\text{tang.}\,\omega = (e^\omega - e^{-\omega})/(e^\omega +
%               e^{-\omega})$ as case II, the solution $y$ with the double
%               sign. Corollarium 1, § 40: $\omega = 5\pi/4, 9\pi/4 \dots$,
%               the same sounds as case II, fundamental $d$ (« vide § 32 »).
%               Corollarium 2, § 41: yet a very different figure, computed at
%               $u = \frac15 \dots \frac45$ (« Tab. III fig. 3 »). Scholion
%               § 42: the general equation with $\Omega$ for
%               $\surd 2(e^{2\omega} + e^{-2\omega})$.
%   v. 187–190  « Evolutio casus VI », both ends built in (Problema § 43):
%               the equation of case I again, $\cos\omega = 2/(e^\omega +
%               e^{-\omega})$, $\omega = 3\pi/2, 5\pi/2 \dots$, the same sounds,
%               only $\gamma$ and $\delta$ of opposite sign. Corollarium I, § 44:
%               the figure of the fundamental computed at $u = \frac13,
%               \frac23$ (« Tab. III fig. 4 », the curve $ErsF$).
%               Corollarium 2, § 45: sounds of similar rods go as $c/aa$,
%               directly as the thickness and inversely as the square of the
%               length. Scholion § 46: rods of rectangular section $ABCD$
%               (« fig. 5 »), $c$ the side in the plane of bending; the
%               circle of diameter $AB$ (« fig. 6 »), $c = \frac34 AB$.
%   v. 192–200  Rods also pinned at an intermediate point $L$ (« fig. 7 »),
%               a case added to show the method (Problema § 47): two
%               portions $EL = \lambda a$ and $LF$, eight coefficients and
%               eight equations, reduced to one transcendental equation in
%               $\omega$ which the text declines to solve; $\lambda = 0$ and
%               $\lambda = 1$ give rest. Then « Evolutio casus » for $L$ at
%               the middle (« Tab. III fig. 8 », § 48): the equation splits
%               into $\sin\frac12\omega$ and a second factor; § 49, the first
%               factor gives case IV for each half, the sounds 4, 16, 36, 64,
%               a double octave above case IV; § 50 the coefficients, which
%               the batch leaves at $\alpha = 0$, $\beta = 0$, $\delta = 0$.
%               The copy runs on into batch 11.
%
% The numbers on the leaves, and why no \folio{} is written. Two kinds, both
% in ink. (1) The copy's parenthesised numbers, one to a copied page, at the
% head of the page in the middle: (140) to (156) at views 181–190, 192–195,
% 198–200, without a gap; at view 189 the last figure of « (14…) » is under
% an ink blot and is not read. As in batch 8 these are very probably the
% pages of the printed volume being copied — an inference. (2) The ink series
% of the earlier batches at the top right of the pages that face up, which
% are the pages with odd copy numbers: 90 at view 182 (141), 91 at 184 (143),
% 92 at 186 (145), 93 at 188 (147), 94 at 190 (149, and again on its second
% photograph, view 191), 95 at 193 (151), 96 at 195 (153, and at view 197),
% 97 at 199 (155). The pages with even copy numbers carry none. No pencilled
% foliation of the Bibliothèque was read anywhere in the batch, so no
% \folio{} is written, as in batches 1 to 8. The marginal « Tab. III fig. 1 »
% to « fig. 8 » are the copy's references to the memoir's plate, which is
% not in these views. Every number above was read on its view; none is
% inferred.
%
% Notation. The copy's letters are kept as in batch 8: $\alpha$, $\beta$,
% $\gamma$, $\delta$ (and primes) for the constants, $\omega = a/f$, $u = s/a$,
% $\zeta$ the phase, $C$ the arbitrary coefficient, $c$, $b$, $g$, $a$ as in
% the memoir; $\Omega$ at view 187; $\lambda$ for the fraction $EL/EF$ from
% view 193, a small hooked letter read as Greek. The radical is written
% without a bar, « $\surd 2gb$ », and set \surd. « sin. », « cos. », « tang. »
% and « cot. » are written with and without their stop in the same line, and
% the pass has not tried to record every stop: where one is set it was seen,
% where it is not the leaf may still carry it. Musical notes above the staff
% are written as letters with one to five bars over them (view 182), set
% with \bar. The copyist's underlining of titles is described in the notes;
% underlined single letters and note names are set \emph{}. Where the
% calculation and the leaf disagree the leaf stands, with a note: view 181
% (a « + » where $\beta = -\alpha$ gives « − »), view 186 (« 56 »), view
% 187 (exponents without $u$, the missing $t$), view 189 (« ipsi ω » for
% $u$), view 194 (two equations numbered III), view 200 ($\beta'$).
%
% What could not be read. Three places are marked \ill{}: the last figure of
% the copy's page number at view 189, under a blot; one or two letters under
% a blot after « fuerit » at view 192; a short word struck with a heavy
% erasure before « sin. $\frac12\omega$ » at view 200. Readings offered with
% doubt are marked \uncertain{}: the sign at view 181; « sequetur »,
% « numerus » and « fere » at view 184; « et » at view 188; « 23,0769 » and
% « 138 » at view 189; « parest » at view 192, which makes no sense and
% probably hides a word under the blot; « poterunt » at view 193; « ausim »
% at view 198; the exponent « 2ω » at view 199. Nothing was guessed to fill
% a gap.
%
% Nothing here was checked against any published transcription or edition
% of Euler's memoir.
\documentclass[11pt,a4paper]{article}
% The preamble shared by every transcription of the mathematicians' manuscripts in Gallica.
%
% It rests on one idea: **uncertainty compounds, it does not resolve.** A
% transcription that smooths over a doubtful word has destroyed the only thing
% it was made for — a reader must be able to tell, at every point, what was on
% the page and what was supplied. The apparatus macros below are therefore the
% critical apparatus, not decoration.
%
% The same commands are recognised by `scripts/render.mjs`, which produces the
% site's reading view, and by `scripts/tei.mjs`. A macro added here must be
% added there too, or rendering fails loudly — which is the wanted behaviour.
%
% Comments here are English, like the rest of the repository. The documents
% this preamble sets are French, and that is deliberate: see the skills.

\ProvidesPackage{germain}[2026/09/15 Transcriptions of mathematicians' manuscripts in Gallica]

\usepackage{amsmath,amssymb,amsthm}
\usepackage{mathrsfs}
% Kept from the parent preamble so the renderer's diagram support stays
% available; these manuscripts draw no commutative diagrams, and nothing here uses it.
\usepackage{tikz-cd}
\usepackage[english,french]{babel}
\usepackage{fontspec}

% --- Greek ------------------------------------------------------------------
% The text face stays fontspec's default, Latin Modern, which has no Greek:
% XeTeX drops every glyph it lacks with a line in the log and nothing on the
% page, and the Greek words the manuscripts quote vanished from the PDFs. So
% the Greek and Greek Extended blocks get a XeTeX character class of their own,
% and entering or leaving a run of it switches the family to CMU Serif — the
% same Computer Modern design, with polytonic Greek — and back. Only the
% family changes: a Greek word in \textit or \textbf keeps its shape and series.
% The transitions to and from every other class carry the tokens that class
% has towards an ordinary letter, so French spacing before « ; » or inside
% « » still holds after a Greek word. No grouping, so no brace or \endgroup
% can come out unbalanced; maths is left alone, since XeTeX inserts nothing
% there. Kept identical in `exercices.sty`.
\makeatletter
\newfontfamily\germainGreekfont{cmun}[
  Extension=.otf, NFSSFamily=germaingreek,
  UprightFont=*rm, ItalicFont=*ti, SlantedFont=*sl,
  BoldFont=*bx, BoldItalicFont=*bi, BoldSlantedFont=*bl]
\newXeTeXintercharclass\germain@greek
\def\germain@greekrange#1#2{%
  \@tempcnta=#1\relax
  \loop
    \XeTeXcharclass\@tempcnta=\germain@greek
    \ifnum\@tempcnta<#2\advance\@tempcnta\@ne\repeat}
\germain@greekrange{"0370}{"03FF}
\germain@greekrange{"1F00}{"1FFF}
\def\germain@greekprev{\rmdefault}
\def\germain@greekin{%
  \edef\germain@greekprev{\f@family}\fontfamily{germaingreek}\selectfont}
\def\germain@greekout{\fontfamily{\germain@greekprev}\selectfont}
\def\germain@greekjoin{%
  \XeTeXinterchartokenstate=\@ne
  \@tempcnta=\z@
  \loop
    \ifnum\@tempcnta=\germain@greek\else
      \XeTeXinterchartoks\germain@greek\@tempcnta=\expandafter{%
        \expandafter\germain@greekout\the\XeTeXinterchartoks\z@\@tempcnta}%
      \XeTeXinterchartoks\@tempcnta\germain@greek=\expandafter{%
        \the\XeTeXinterchartoks\@tempcnta\z@\germain@greekin}%
    \fi
    \ifnum\@tempcnta<4095\advance\@tempcnta\@ne\repeat}
% babel-french sets its own classes, and saves and restores the interchar
% state, each time a language is selected: join again after it does.
\addto\extrasfrench{\germain@greekjoin}
\addto\extrasenglish{\germain@greekjoin}
\AtBeginDocument{\germain@greekjoin}
\makeatother

\usepackage{geometry}
\geometry{a4paper,margin=25mm,marginparwidth=28mm}
\usepackage{marginnote}
\usepackage{xcolor}
\usepackage[normalem]{ulem}
\setlength{\emergencystretch}{3em}
\usepackage{hyperref}
\hypersetup{colorlinks=true,linkcolor=black,urlcolor=black,pdfborder={0 0 0}}

% --- Batch metadata ---------------------------------------------------------
% Taken as they stand by the rendering script, which shows them at the head of
% the reading view. They fix what the file claims to cover. The macro names are
% the parent project's, so that the scripts stay shared; \folder here means the
% volume's slug — `fr-9115` — and \pages counts Gallica views.

\newcommand{\folder}[1]{\gdef\grothFolder{#1}}
\newcommand{\batch}[1]{\gdef\grothBatch{#1}}
\newcommand{\pages}[2]{\gdef\grothFirst{#1}\gdef\grothLast{#2}}
\newcommand{\foldertitle}[1]{\gdef\grothFoldertitle{#1}}

% The holder's own shelfmark, printed as they write it: « Français 9115 ».
\newcommand{\shelfmark}[1]{\gdef\grothShelfmark{#1}}

% The dating, copied verbatim from the holder's catalogue — never inferred
% here. « XIXe siècle » is the BnF's; a pass that dates a leaf from its content
% says so in a \note{}, as its own inference.
\newcommand{\dating}[1]{\gdef\grothDating{#1}}

% The status notice, printed in the title block and repeated as a diagonal
% watermark on every page of the PDF. The manuscripts are in the public
% domain; what this line declares is not a right but a quality — an
% unverified first pass by a machine. The skills write the line; a file
% without it gets no watermark, so leaving it out is visible in the output.
\newcommand{\watermark}[1]{\gdef\grothWatermark{#1}}

\gdef\grothFolder{?}\gdef\grothBatch{}\gdef\grothFirst{?}\gdef\grothLast{?}%
\gdef\grothFoldertitle{}\gdef\grothShelfmark{}\gdef\grothDating{}\gdef\grothWatermark{}

% --- The résumé -------------------------------------------------------------
\newenvironment{resume}
  {\small\setlength{\parskip}{0.45em}}
  {\par\medskip\hrule\medskip}

% --- Keywords ---------------------------------------------------------------
% In English, closing the modernised reading's résumé: the modern vocabulary
% under which to search for what the leaves construct. The manifest extracts
% them to tag the volume — this line is the single source of the tags.
\newcommand{\keywords}[1]{\par\smallskip\noindent
  {\footnotesize\textsc{Keywords} --- \textit{#1}}\par}

% --- The critical apparatus -------------------------------------------------

% The Gallica view number, in the margin. This is the anchor that lets a
% disagreement over a formula be settled by looking at *one* image rather than
% a volume of 750. The site uses it to turn the facsimile as the transcription
% is scrolled. A view is one scanned image: usually one side of a leaf.
\newcommand{\page}[1]{%
  \marginnote{\footnotesize\color{gray!70}\textsf{v.\,#1}}%
  \label{page:#1}\ignorespaces}

% The BnF's pencilled foliation, where the leaf carries it and a pass can read
% it: « 198r ». This is what the literature cites — « ff. 198r–208v » — and
% the one line that turns a citation into a view. Recorded, never inferred:
% a folio a pass cannot read is not written.
\newcommand{\folio}[1]{%
  \marginnote{\footnotesize\color{gray!50}\textsf{f.\,#1}}\ignorespaces}

% The modernised reading's coarser counterpart to \page{}: which views a
% section is drawn from.
\newcommand{\pagerange}[2]{%
  \marginnote{\footnotesize\color{gray!70}\textsf{v.\,#1--#2}}%
  \ignorespaces}

% Illegible: never guessed.
\newcommand{\ill}{\textcolor{red!60!black}{[\,\ldots\,]}}

% A doubtful reading: the word is given, and flagged as uncertain.
\newcommand{\uncertain}[1]{\uline{#1}}

% An editorial addition — a missing letter, an implied word.
\newcommand{\add}[1]{\textcolor{blue!50!black}{[#1]}}

% Struck out by the writer. What was crossed out is part of the document.
\newcommand{\struck}[1]{\sout{#1}}

% The transcriber's note, on the state of the page or the mathematics.
\newcommand{\note}[1]{\par\smallskip{\small\color{gray!60!black}\textit{[#1]}}\par\smallskip}

% A marginal note of hers, distinct from the body of the text.
\newcommand{\marginal}[1]{\par\smallskip\noindent\hspace{1em}%
  {\small\color{gray!60!black}#1}\par\smallskip}

% --- Title ------------------------------------------------------------------
% The title block is French, like the documents themselves.

\AddToHook{shipout/background}{%
  \ifx\grothWatermark\empty\else
    \begin{tikzpicture}[remember picture, overlay]
      \node[rotate=52, scale=3.4, align=center, text=gray!18,
            font=\sffamily\bfseries]
        at (current page.center) {\grothWatermark};
    \end{tikzpicture}%
  \fi}

\AtBeginDocument{%
  \begin{center}
    {\large\bfseries Manuscrits de mathématiciens (Gallica) ---
      \ifx\grothShelfmark\empty\grothFolder\else\grothShelfmark\fi}\\[2pt]
    {\itshape\grothFoldertitle}\\[4pt]
    {\small vues \grothFirst--\grothLast
      \ifx\grothBatch\empty\else\ (lot \grothBatch)\fi}\\[2pt]
    \ifx\grothDating\empty\else
      {\small\color{gray!60!black}Datation du catalogue : \grothDating}\\[2pt]
    \fi
    \ifx\grothWatermark\empty\else
      {\small\bfseries\color{red!45!gray}\grothWatermark}\\[2pt]
    \fi
    {\footnotesize\color{gray!60!black}Transcription établie d'après les images IIIF de Gallica.
     Source gallica.bnf.fr / Bibliothèque nationale de France.\\
     Les lectures incertaines sont soulignées, les illisibles notées [\,\ldots\,],
     les ajouts entre crochets ; \textsf{v.} numérote les vues de Gallica,
     \textsf{f.} la foliotation au crayon de la BnF.}
  \end{center}
  \vspace{1em}}

\endinput


\folder{fr-9115}
\shelfmark{Français 9115}
\batch{10}
\pages{181}{200}
\dating{1801-1900}
\watermark{Lecture automatique\\non vérifiée}
\foldertitle{Papiers de Sophie GERMAIN. — Recueil de dissertations et problèmes mathématiques et physiques.}

\begin{document}

\note{Numérotation des feuillets. La copie latine porte en tête de chaque
page, au milieu, un nombre entre parenthèses, de (140) à (156), qui semble
reporter la page de l'imprimé copié : c'est une inférence ; à la vue 189, le
dernier chiffre de « (14…) » est sous une tache d'encre et n'est pas lu. La
série à l'encre des lots précédents se poursuit en haut à droite des pages
de numéro impair : 90 (vue 182), 91 (184), 92 (186), 93 (188), 94 (190), 95
(193), 96 (195), 97 (199). Tous ces nombres ont été lus sur la vue ; aucun
n'est déduit. Aucune foliotation au crayon de la Bibliothèque n'a été lue
sur ces vues, et aucun « f. » n'est donc porté ici. Les vues 191, 196 et 197
photographient une seconde fois les pages des vues 190, 194 et 195 et ne
sont pas transcrites.}

\page{181}
(140)

\section{Evolutio casus IV. quo virgae uterque terminus simpliciter stylo est fixus.}

\subsection{Problema}

§. 35. Si virgae elasticae tam terminus $E$ quam $F$ simpliciter fuerit
fixus, investigare omnes motus, quibus ea contremiscere potest.

\emph{Solutio.}

Posito igitur tam $s=0$ quam $s=a$, utroque casu fieri oportet et $y=0$ et
$\left(\frac{ddy}{ds^2}\right)=0$; ex priore casu hae deducuntur aequationes:
\[ 1^{\circ})\ \alpha+\beta+\delta=0, \quad\text{et}\quad 2^{\circ})\ \alpha+\beta-\delta=0, \]
ex quibus statim colligitur
\[ \alpha+\beta=0 \text{ et } \delta=0 \text{ ideoque } \beta=-\alpha. \]
Alter vero casus, quo $s=a$, ponendo $\frac{a}{f}=\omega$, has suppeditat
aequationes:
\[ 3^{\circ})\ \alpha e^{\omega}+\beta e^{-\omega}+\gamma \text{sin.}\,\omega+\delta\cos.\omega=0 \]
\[ 4^{\circ})\ \alpha e^{\omega}+\beta e^{-\omega}-\gamma \sin \omega-\delta\cos.\omega=0, \]
quae superioribus valoribus substitutis, reducuntur ad has:
\[ \alpha(e^{\omega}\uncertain{+}e^{-\omega})+\gamma\text{sin.}\,\omega=0 \text{ et } \alpha(e^{\omega}-e^{-\omega})-\gamma\text{sin.}\,\omega=0, \]
\note{le signe entre $e^{\omega}$ et $e^{-\omega}$, dans la première de ces deux équations, a la forme d'un « + » ; la substitution $\beta=-\alpha$ donne un « − », et la somme qui suit est écrite avec « − ». Laissé tel que lu.}
quarum summa praebet
\[ 2\alpha(e^{\omega}-e^{-\omega})=0 \text{ ideoque } \alpha=0, \]
differentia vero dat $2\gamma\text{sin.}\,\omega=0$, unde si sumeremus $\gamma=0$, tota
\note{« (140) » en haut, au milieu. « Evolutio casus » est souligné, et la seconde ligne du titre l'est d'un bout à l'autre ; « Problema » et « Solutio » sont soulignés. Le feuillet s'arrête sur « tota », dont la phrase se poursuit à la vue 182.}

\page{182}
(141)

virga in quiete esset mansura; necesse igitur est ut sit $\sin \omega=0$. Hinc ergo
statim innotescunt omnes valores pro angulo $\omega$, quippe qui sunt:
\[ 0=\omega,\ \omega=\pi,\ \omega=2\pi,\ \omega=3\pi,\ \omega=4\pi\ \&c. \]
quorum primus locum habet in statu quietis, secundus autem sonum
fundamentalem hoc casu exhibet, ex quo, ut in praecedentibus casibus,
oriuntur sequentes soni:
\[ \frac{\pi c\sqrt{2gb}}{aa},\ \frac{4\pi c\sqrt{2gb}}{aa},\ \frac{9\pi c\sqrt{2gb}}{aa},\ \frac{16\pi c\sqrt{2gb}}{aa}\ \&c. \]
qui igitur soni secundum numeros quadratos $1, 4, 9, 16$ ascendunt, ita ut,
si fundamentalis fuerit $C$, isti soni constituant hanc seriem:
$C$, $\bar{c}$, $\bar{\bar{d}}$, $\bar{\bar{\bar{c}}}$, $\bar{\bar{\bar{\bar{\text{gis}}}}}$, $\bar{\bar{\bar{\bar{\bar{d}}}}}$ \&c. quod si ergo omnes istos motus generaliter
conjungamus, formula generalis, omnes plane motus, quos virga
nostra recipere potest complectens, erit
\begin{align*}
y ={}& C\sin \Big(\zeta+t\frac{\pi\pi c\sqrt{2gb}}{aa}\Big)\text{sin.}\,\pi u + C'\sin \Big(\zeta'+t\frac{4\pi\pi c\sqrt{2gb}}{aa}\Big)\text{sin.}\,2\pi u \\
&+ C''\sin \Big(\zeta''+t\frac{9\pi\pi c\sqrt{2gb}}{aa}\Big)\text{sin.}\,3\pi u + C'''\sin \Big(\zeta'''+t\frac{16\pi\pi c\sqrt{2gb}}{aa}\Big)\sin 4.\pi u \\
&+ \&c.
\end{align*}
ubi scribamus $u$ loco $\frac{s}{a}$.
\note{notes de musique : les barres au-dessus des lettres sont au nombre de une ($c$), deux ($d$), trois ($c$), quatre ($gis$) et cinq ($d$), telles qu'on les compte sur la vue ; « gis » est de plus souligné. Entre « igitur » et « est », un trait vertical à la plume. Dans les sons, le numérateur porte un seul $\pi$ ; dans la formule générale, deux. « \&c. » est écrit deux fois à la fin de la formule, à la suite du dernier terme et à la ligne.}

\subsection{Corollarium 1.}

\marginal{Tab. III. fig. 1.}
§. 36. Quando ergo virga sonum edit fundamentalem, eandem
recipiet curvaturam quam chordae simpliciter vibrantes, quippe
quae erit linea sinuum, qualem figura prima refert. Pro sono autem
simplice secundo, quo $\omega=2\pi$, figura chordae erit (fig. 2), habens unum nodum
in medio $O$. Pro sequentibus autem sonis simplicibus numerus nodorum semper
unitate augetur.
\note{« (141) » en haut, au milieu ; le nombre 90 est à l'encre en haut à droite. « Corollarium 1. » est souligné.}

\page{183}
(142)

§. 37. Quod si sonos fundamentales omnium horum quatuor casuum,
quos hactenus tractavimus, inter se comparemus, jam vidimus, si \struck{sonus}
sonus primi casus exprimatur per \emph{gis}, pro secundo casu eum fore $d$,
\struck{\&c} pro tertio \emph{C}, qui soni his numeris exprimuntur: $\frac{9}{4}$, $\frac{25}{16}$, $\frac{9}{25}$.
Hinc cum praesenti casu quarto sonus fundamentalis exprimatur
unitate, sonus erit \emph{fis}; seriem igitur hoc modo referamus:
\[ \text{I.}\ \begin{array}{c}\frac{9}{4}\\ \text{gis}\end{array}.\quad
\text{II.}\ \begin{array}{c}\frac{25}{16}\\ d\end{array}:\quad
\text{III.}\ \begin{array}{c}\frac{9}{25}\\ C\end{array}.\quad
\text{IV}\ \begin{array}{c}1.\\ \text{fis}\end{array} \]

\subsection{Scholion}

§. 38. Hic igitur casus, quo uterque virgae terminus simpliciter
est fixus, prae reliquis hac insigni gaudet praerogativa, quod
omnes valores anguli $\omega$ accurate sine ullo errore definire liceat,
propterea quod formulae exponentiales, quae hanc determinationem
turbabant et non parum irregularem reddebant, penitus ex calculo
evanuerunt; unde soni hoc casu editi multo magis ad harmoniam
sunt \struck{accomodati} accommodati, et quidem adhuc magis quam in
chordis simplicibus usu venit. Cum enim soni ab eadem chorda
editi secundum numeros $1, 2, 3, 4$ \&c. progrediantur, fere semper
plures horum sonorum simul audiuntur, inter quos etiam non
parum dissoni occurrere possunt. Verum quia a nostra virga
alii soni edi non possunt, nisi qui numeris $1, 4, 9, 16, 25$ exprimuntur,
praeter fundamentalem potissimum exaudietur ejus dupla octava
\note{les noms de notes sont soulignés sur le feuillet (\emph{gis}, \emph{C}, \emph{fis}) ; le « C » du tableau l'est aussi. Entre « praeter » et « fundamentalem », un trait vertical à la plume, comme à la vue 182.}
\note{« (142) » en haut, au milieu. « Scholion » est souligné. Le feuillet s'arrête sur « octava », la phrase se poursuit à la vue 184.}

\page{184}
(143)

harmoniam nihil turbans, postea vero \uncertain{sequetur} sonus \uncertain{numerus} 9 respondens,
qui, cum fundamentalem ultra \uncertain{fere} tres octavas superet, ob nimium acumen
vix unquam audietur, ita ut tantum dupla octava simul cum fundamentali
tinniat. Talis igitur virga sonos multo puriores reddere est censenda
quam chordae simplices, unde etiam soni hoc modo editi in musica
peculiarem suavitatem habere debebunt.
\note{« sequetur » est écrit en deux morceaux séparés par un blanc, et la lettre du milieu n'est pas nette ; la fin de « numerus » a la forme de « -us » ; « fere » a une première lettre sans hampe, qui ressemble à un « t ».}

\section{Evolutio casus V quo virgae elasticae terminus simpliciter est fixus, alter vero quasi muro firmiter infixus}

\subsection{Problema.}

§. 39. Si virgae elasticae terminus $E$ fuerit simpliciter fixus,
alter vero $F$ prorsus infixus, investigare omnes motus, quibus ea
contremiscere potest.

\emph{Solutio.}

Quia terminus $E$, ubi fit $s=0$, simpliciter est fixus, habebimus
statim, uti in casu praecedente, $\beta=-\alpha$ et $\delta=0$; pro altero autem
termino, ubi $s=a$, erit tam $y=0$ quam $\left(\frac{dy}{ds}\right)=0$, unde posito $\frac{a}{f}=\omega$
oriuntur hae duae aequationes:
\[ \alpha e^{\omega}+\beta e^{-\omega}+\gamma\text{sin.}\,\omega+\delta\cos.\omega=0, \]
\[ \alpha e^{\omega}-\beta e^{-\omega}+\gamma\cos.\omega-\delta\text{sin.}\,\omega=0, \]
quae, substitutis praecedentibus valoribus, abeunt in has:
\[ \alpha(e^{\omega}-e^{-\omega})+\gamma\sin \omega=0 \quad\text{et} \]
\[ \alpha(e^{\omega}+e^{-\omega})+\gamma\cos.\omega=0. \]
\note{une tache d'encre devant « fixus », à la première ligne du §. 39.}
\note{« (143) » en haut, au milieu ; le nombre 91 est à l'encre en haut à droite. Les deux lignes du titre sont soulignées chacune d'un trait, « Problema » et « Solutio » aussi.}

\page{185}
(144)

ex priore fit
\[ \frac{\alpha}{\gamma}=\frac{-\sin \omega}{e^{\omega}-e^{-\omega}}, \text{ ex altero vero } \frac{\alpha}{\gamma}=\frac{-\cos\omega}{e^{\omega}+e^{-\omega}}; \]
ex quibus porro colligitur $\text{tang.}\,\omega=\frac{e^{\omega}-e^{-\omega}}{e^{\omega}+e^{-\omega}}$, quae formula plane
convenit cum ea, quam casu secundo invenimus, eritque idcirco,
ut ibi, vel
\[ \text{sin.}\,\omega=\frac{e^{\omega}-e^{-\omega}}{\surd 2(e^{2\omega}+e^{-2\omega})} \text{ et } \cos.\omega=\frac{e^{\omega}+e^{-\omega}}{\surd 2(e^{2\omega}+e^{-2\omega})}, \text{ vel} \]
\[ \text{sin.}\,\omega=\frac{-e^{\omega}+e^{-\omega}}{\surd 2(e^{2\omega}+e^{-2\omega})} \text{ et } \cos.\omega=\frac{-e^{\omega}-e^{-\omega}}{\surd 2(e^{2\omega}+e^{-2\omega})}. \]
Ex prioribus valoribus elicitur $\frac{\alpha}{\gamma}=\frac{-1}{\surd 2(e^{2\omega}+e^{-2\omega})}$, ex alteris
autem valoribus fit $\frac{\alpha}{\gamma}=\frac{+1}{\surd 2(e^{2\omega}+e^{-2\omega})}$; unde his binis
casibus conjungendis habebimus $\alpha=1$, $\beta=-1$, \struck{$\gamma=\surd 2($}
$\gamma=\mp\surd 2(e^{2\omega}+e^{-2\omega})$ et $\delta=0$; ubi ut hactenus signorum
ambiguorum superius valet, si anguli $\omega$ tam sinus quam
cosinus fuerint positivi, inferius autem si ambo fuerint
negativi, unde pro quovis valore $\omega$ habebitur
\[ y=C\sin \Big(\zeta+t\frac{\surd 2g}{k}\Big)\big\{e^{u\omega}-e^{-u\omega}\mp\surd 2(e^{2\omega}+e^{-2\omega})\text{sin.}\,u\omega\big\}; \]
praeterea vero ut hactenus erit $\frac{\surd 2g}{k}=\frac{\omega\omega c\surd 2gb}{aa}$, tempus
unius oscillationis $=\frac{\pi aa}{\omega\omega c\surd 2gb}$ et sonus editus $=\frac{\omega\omega c\surd 2gb}{\pi aa}$.

\subsection{Corollarium 1.}

§. 40. Quia aequatio resolvenda: $\text{tang.}\,\omega=\frac{e^{\omega}-e^{-\omega}}{e^{\omega}+e^{-\omega}}$
prorsus convenit cum ea, quam casu secundo jam resolvimus
\note{« (144) » en haut, au milieu. Le radical est écrit sans barre, « $\surd 2(e^{2\omega}+e^{-2\omega})$ », et il est transcrit ainsi. Le « $\gamma=\surd 2($ » biffé en fin de ligne est barré d'un trait, le $\gamma$ d'une croix. Le double signe de $\gamma$ et de la formule de $y$ est un « − » posé sur un « + ». « Corollarium 1. » est souligné. Le feuillet s'arrête sur « resolvimus ».}

\page{186}
(145)

omnes valores anguli $\omega$ erunt, ut ibi, sequentes:
\[ \omega=\frac{5\pi}{4},\ \frac{9\pi}{4},\ \frac{13\pi}{4},\ \frac{17\pi}{4},\ \&c. \]
unde etiam omnes soni simplices, quos haec virga edere potest,
sequentibus numeris exprimuntur:
\[ \frac{25\pi c\surd 2gb}{16aa};\ \frac{81\pi c\surd 2gb}{16aa};\ \frac{169.\pi c\surd 2gb}{16aa};\ \frac{289\pi c\surd 2gb}{16aa};\ \&c. \]
quorum primus etiam pro fundamentali habetur, et secundum
superiorem determinationem respondet clavi $d$ (vide §. 32)

\subsection{Corollarium 2.}

§. 41. Quanquam autem hic casus eosdem plane sonos simplices
producit, quos casu secundo invenimus, tamen ipsa virga maxime
diversas recipit figuras. Ita pro sono fundamentali, ubi $\omega=\frac{5\pi}{4}$,
ideoque tam sinus quam cosinus sunt negativi, scilicet $\sin \omega=\cos\omega
=-\frac{1}{\surd 2}$, omisso coefficiente erit
\[ y=\text{- - -}\ e^{u\omega}-e^{-u\omega}+\surd 2(e^{2\omega}+e^{-2\omega})\sin u\omega, \]
qui valor pro utroque termino $u=0$ et $u=1$ fit $=0$.
\note{entre « $y=$ » et « $e^{u\omega}$ », trois petits traits en ligne, comme des points de conduite ; ils sont reproduits tels quels.}
Pro reliqua figura cognoscenda, quia supra jam vidimus
esse $e^{2\omega}=e^{\frac{5}{2}\pi}=2576$, erit $\surd 2(e^{2\omega}+e^{-2\omega})=72$ proxime. Tribuamus
igitur literae $u$ sequentes valores: $\frac{1}{5}$, $\frac{2}{5}$, $\frac{3}{5}$, $\frac{4}{5}$, et pro singulis
sequentes valores ipsius $y$ prodibunt:

\marginal{Tab. III fig. 3}
\begin{enumerate}
\item[I.] Sit $u=\frac{1}{5}$, erit $u\omega=\frac{\pi}{4}=45^{\circ}$ et $y=1,7974+\frac{72}{\surd 2}=56$ proxime;
\item[II.] Sit $u=\frac{2}{5}$ erit $u\omega=\frac{\pi}{2}=90^{\circ}$ et $y=4,6085+72=76$ proxime;
\item[III.] Si $u=\frac{3}{5}$ erit $u\omega=\frac{3\pi}{4}=135^{\circ}$ et $y=10,46+\frac{72}{\surd 2}=61$ proxime;
\end{enumerate}
\note{à la ligne I, la somme $1,7974+\frac{72}{\surd 2}$ vaut à peu près 52,7 et non 56 ; le chiffre lu est « 56 », et il est laissé tel. Le « 3 » de $\frac{3}{5}$ (ligne III) est tracé d'une boucle qui le rapproche d'un « 9 ».}
\note{« (145) » en haut, au milieu ; le nombre 92 est à l'encre en haut à droite. « Corollarium 2. » est souligné, et la lettre $u$ de « literae u ». Le feuillet s'arrête après la ligne III ; la suite de la table est à la vue 187.}

\page{187}
(146)

\begin{enumerate}
\item[IV.] Sit $u=\frac{4}{5}$, erit $u\omega=\pi=180^{\circ}$ et $y=23,10+0=23$ proxime.
\end{enumerate}

Ista figura repraesentatur per figuram tertiam;

\subsection{Scholion}

§. 42. Ut jam omnes plane motus ex inventis simplicibus
concinne repraesentemus, pro formula irrationali $\surd 2(e^{2\omega}+e^{-2\omega})$
scribamus characterem $\Omega$, cui pro variis valoribus $\omega, \omega', \omega'', \omega'''$ \&c.
\struck{tribu} tribuamus valores $\Omega, \Omega', \Omega'', \Omega'''$ \&c et scribendo, ut hactenus,
$u$ loco $\frac{s}{a}$, aequatio generalis erit:
\begin{align*}
y ={}& C\,\text{sin.}\Big(\zeta+\frac{\omega\omega c\surd 2gb}{aa}t\Big)(e^{\omega}-e^{-\omega}\mp\Omega\,\text{sin.}\,u\omega) \\
&+ C'\,\text{sin.}\Big(\zeta'+\frac{\omega'\omega'c\surd 2gb}{aa}t\Big)(e^{\omega'}-e^{-\omega'}\mp\Omega'\,\text{sin.}\,u\omega') \\
&+ C''\,\text{sin.}\Big(\zeta''+\frac{\omega''\omega''c\surd 2gb}{aa}\Big)(e^{\omega''}-e^{-\omega''}\mp\Omega''\sin u\omega'') \\
&+ C'''\,\text{sin.}\Big(\zeta'''+\frac{\omega'''\omega'''c\surd 2gb}{aa}\Big)(e^{\omega'''}-e^{-\omega'''}\mp\Omega'''\,\text{sin.}\,u\omega''')
\end{align*}
ubi signum superius valet pro angulis $\omega, \omega', \omega''$ \&c. quorum
sinus et cosinus sunt positivi, inferius autem ubi sunt negativi.

\section{Evolutio casus VI. quo virgae elasticae uterque terminus firmiter quasi muro est infixus}

\subsection{Problema}

§. 43 Si virgae elasticae uterque terminus firmiter fuerit
infixus, investigare motus, quibus ea contremiscere potest

\emph{Solutio}

Hoc igitur casu pro utroque termino debet esse tam $y=0$ quam $\left(\frac{dy}{ds}\right)=0$;
\note{« (146) » en haut, au milieu. La formule générale est transcrite telle qu'elle est écrite : les exposants des exponentielles ne portent pas le facteur $u$ que porte le sinus ; le $t$ suit la fraction dans les deux premières lignes et manque dans les deux dernières ; la formule n'est pas close par « \&c ». Le « tribu » biffé est au début de la ligne. « Scholion », « Evolutio casus VI. », la seconde ligne du titre, « Problema » et « Solutio » sont soulignés. Le feuillet s'arrête sur « $=0$; » ; le bas de la page est blanc.}

\page{188}
(147)

pro priore termino ergo habebimus has aequationes:
\[ \alpha+\beta+\delta=0 \text{ et } \alpha-\beta+\gamma=0, \]
unde consequimur $\delta=-\alpha-\beta$ et $\gamma=-\alpha+\beta$; posterior vero, quo $s=a$
et $\frac{a}{f}=\omega$, praebet
\[ \alpha e^{\omega}+\beta e^{-\omega}+\gamma\,\text{sin.}\,\omega+\delta\cos.\omega=0 \text{ et} \]
\[ \alpha e^{\omega}-\beta e^{-\omega}+\gamma\cos.\omega-\delta\,\text{sin.}\,\omega=0; \]
ubi si priores valores substituantur, oriuntur hae equationes:
\[ \alpha(e^{\omega}-\text{sin.}\,\omega-\cos.\omega)+\beta(e^{-\omega}+\text{sin.}\,\omega-\cos.\omega)=0 \text{ et} \]
\[ \alpha(e^{\omega}-\cos.\omega+\text{sin.}\,\omega)-\beta(e^{-\omega}-\cos\omega-\text{sin.}\,\omega)=0; \]
unde duplici modo colligitur
\[ \frac{\alpha}{\beta}=\frac{-e^{-\omega}-\sin\omega+\cos\omega}{e^{\omega}-\sin\omega-\cos.\omega}=\frac{e^{-\omega}-\text{sin.}\,\omega-\cos.\omega}{e^{\omega}+\text{sin.}\,\omega-\cos.\omega}, \]
qui valores cum prorsus conveniat cum iis, quos casu primo
sumus nacti, inde etiam sequitur fore $\cos.\omega=\frac{2}{e^{\omega}+e^{-\omega}}$; sicque
etiam omnes valores anguli $\omega$ iidem erunt, qui primo casu
jam sunt exhibiti scilicet $\omega=\frac{3\pi}{2}, \frac{5\pi}{2}, \frac{7\pi}{2}$ \&c. unde etiam
ista virga eosdem edet sonos, qui erant:
\[ \frac{9\pi c\surd 2gb}{4aa};\ \frac{25\pi c\surd 2gb}{4aa};\ \frac{49\pi c\surd 2gb}{4aa};\ \frac{81\pi c\surd 2gb}{4aa};\ \&c. \]
\struck{Praete} Praeterea vero etiam erit $\text{sin.}\,\omega=\pm\frac{(e^{\omega}-e^{-\omega})}{e^{\omega}+e^{-\omega}}$, ubi
signum superius valet pro casibus ubi sinus est positivus,
inferius si negativus; unde porro obtinemus $\alpha=1$, $\beta=\mp e^{\omega}$, hincque
porro $\gamma=-(1\pm e^{\omega})$ \uncertain{et} $\delta=-(1\mp e^{\omega})$. Aequatio igitur pro motu a casu
primo in eo tantum differt, quod hic coefficientes $\gamma$ et $\delta$ contraria
\note{« (147) » en haut, au milieu ; le nombre 93 est à l'encre en haut à droite. Entre « habebimus » et « has », un trait vertical à la plume. Le mot biffé au début de l'avant-dernier alinéa est un « Praete » inachevé, repris aussitôt en « Praeterea ». Entre les deux valeurs de $\gamma$ et de $\delta$, le mot de liaison est réduit à un signe serré qui ressemble à un « + » ; il est lu « et » avec doute. Le « a » de « a casu » est ajouté au-dessus de la ligne avec un signe d'insertion ; il est mis à sa place. Trois taches d'encre au milieu de la ligne « signum superius valet pro casibus ». Le feuillet s'arrête sur « contraria ».}

\page{189}
(14\ill{})

signa sunt nacti; sicque, si $\omega, \omega', \omega'', \omega'''$ denotent omnes valores ipsius $\omega$,
aequatio generalis, omnes motus virgae complectens, erit:
\begin{align*}
y ={}& C\,\text{sin.}\Big(\zeta+t.\frac{\omega\omega c\surd 2gb}{aa}\Big)\big\{e^{u\omega}\mp e^{\omega(1-u)}-(1\pm e^{\omega})\,\text{sin.}\,u\omega-(1\mp e^{\omega})\cos.u\omega\big\} \\
&+ C'\,\text{sin.}\Big(\zeta'+t.\frac{\omega'\omega'c\surd 2gb}{aa}\Big)\big\{e^{u\omega'}\mp e^{\omega'(1-u)}-(1\pm e^{\omega'})\,\text{sin.}\,u\omega'-(1\mp e^{\omega'})\cos.u\omega'\big\} \\
&+ C''\,\text{sin.}\Big(\zeta''+t.\frac{\omega''\omega''c\surd 2gb}{aa}\Big)\big\{e^{u\omega''}\mp e^{\omega''(1-u)}-(1\pm e^{\omega''})\,\text{sin.}\,u\omega''-(1\mp e^{\omega''})\cos.u\omega''\big\} \\
&\&c.\quad \&c \quad \&c.
\end{align*}

\subsection{Corollarium I.}

§. 44. Quanquam autem omnes motus cum casu primo
perfecte conveniunt: tamen figura, quam virga inter
vibrandum recipit, toto coelo est diversa. Ad quod \struck{ostendam}
ostendendum evoluamus figuram pro sono fundamentali,
ubi est $\omega=\frac{3\pi}{2}$, cujus sinus cum sit negativus, signa valebunt
inferiora, eritque omisso coefficiente
\[ y=\ldots\ e^{u\omega}+e^{\omega(1-u)}-(1-e^{\omega})\,\text{sin.}\,u\omega-(1+e^{\omega})\cos.u\omega, \]
ubi est $e^{\omega}=111$. Nunc autem ipsi $\omega$ tribuamus duos valores
$\frac{1}{3}$ et $\frac{2}{3}$, ac reperiemus:

\marginal{Tab. III fig. 4}
\begin{enumerate}
\item[I.] Si $u=\frac{1}{3}$, $u\omega=\frac{\pi}{2}=90^{\circ}$ et $e^{u\omega}=4,8104$ et $e^{-u\omega}=0,2079$
hincque $y=4,8104+\uncertain{23,0769}+110=\uncertain{138}$.
\item[II] Si $u=\frac{2}{3}$, $u\omega=\pi=180^{\circ}$ et $e^{u\omega}=23,140$ et $e^{\omega(1-u)}=4,810$,
hincque $y=23,140+4,810+112=140$. Si autem duo valores, si
calculum accuratius instituissemus, prodiissent aequales. Curva igitur,
\end{enumerate}
\note{« (14… ) » en haut, au milieu : le dernier chiffre est noyé sous une tache d'encre dont ne dépasse qu'une boucle, et n'est pas lu. La page n'a pas de nombre en haut à droite. Entre « nacti; » et « sicque », un trait vertical à la plume. Dans la formule de $y$ du Corollaire, une ligne de points suit le signe d'égalité, comme à la vue 186. « ipsi ω tribuamus » : la lettre est écrite comme l'$\omega$ de la page, bien que les valeurs $\frac{1}{3}$ et $\frac{2}{3}$ soient ensuite données à $u$. À la ligne I, les chiffres de la somme sont surchargés : le premier « 4 » est repris sur un autre chiffre, et le second terme, lu « 23,0769 », est écrit sur une première écriture qui le rend douteux ; le total, lu « 138 », a un deuxième chiffre dont la boucle le rapproche d'un « 9 ». « Corollarium I. » est souligné. Le feuillet s'arrête sur « Curva igitur, ».}

\page{190}
(149)

quam virga hoc casu induit, habebit figuram $ErsF$, fig. 4 repraesentatam.
Pro sequentibus autem sonis simplicibus vel unus nodus, vel duo, vel tres
successive ingredientur

\subsection{Corollarium 2}

§. 45. Ope formularum, quas pro singulis his casibus eruimus, non
solum omnes soni, quos eadem virga elastica, diversimode constituta,
edere valet; inter se comparari possunt, sed etiam soni diversarum
virgarum, quae tum longitudine quam crassitie inter se discrepant, dijudicari possunt, dummodo crassities in omnibus fuerit
similis: veluti si virgae fuerint cylindricae, quo casu crassities
cujusque circulo repraesentatur; si enim talium virgarum diameter
crassitiei fuerit $=c$, longitudo $=a$, tum sub similibus circumstantiis
soni erunt ut $\frac{c}{aa}$, hoc est directe ut diameter crassitiei et
reciproce ut quadratum longitudinis, ita ut quo crassior fuerit virga
pro eadem longitudine, sonus edatur tanto acutior; contra vero
quo longior fuerit virga pro eadem crassitie, eo gravior sonus
sit proditurus, idque in ratione duplicata
\note{« crassitie inter se discrepant, dijudicari possunt, dummodo » est écrit au-dessus de la ligne, entre un « + » placé après « quam » et un astérisque placé sous « crassities », avec une croix en tête ; l'addition est mise à sa place.}

\subsection{Scholion}

§. 46. Hic scilicet assumsimus, in diversis virgis crassitiem similem
figuram habere, veluti circularem, quippe quo casu virgae sunt cylindricae
et versus omnes plagas aequaliter inflexioni resistunt. Verum etiam
nostrae formulae ad ejusmodi virgas applicari possunt quarum crassities
alia quacunque figura exhibebant. Ad quod ostendendum consideremus ejusmodi
virgam, cujus sectiones transversales ad longitudinem normaliter
\note{« (149) » en haut, au milieu ; le nombre 94 est à l'encre en haut à droite. Entre « sequentibus » et « autem », un trait vertical à la plume. « Corollarium 2 » et « Scholion » sont soulignés. Le feuillet s'arrête sur « normaliter ».}

\page{192}
(150)

\marginal{Tab. III fig. 5}
factae sint parallelogramma rectangula $ABCD$, in quibus ergo duplex
potissimum inflexio locum habere potest, quarum altera fit secundum axem
$AB$, quando scilicet linea $AC$ et $BD$ circa hunc axem inflectuntur; altera
autem inflexio principalis fieri potest circa axem $AC$. Illo casu litera
nostra \emph{c}, quae superioribus formulis inest, aequabitur lateri $AC$,
posteriore vero casu lateri $AB$, utroque vero casu alterum latus plane
non in computum venitur. Litera enim \emph{b} pendet, uti jam supra
notavimus, ab elasticitate absoluta materiae, ex qua virgae sunt
fabricatae. Ita si virga cujus longitudo est $=a$, incurvetur, circa
axem $AB$, sonus editus erit cum $\frac{AC}{aa}$; et si virga incurvetur circa
axem $AC$ sonus editus erit ut $\frac{AB}{aa}$, si scilicet reliquae circumstantiae
fuerint pares. At si sectio virgae transversalis fuerit \ill{} \uncertain{parest}.
Circulus, diametro $AB$ descriptus (fig. 6) tum pro formulis nostris
erit $c=\frac{3}{4}AB$; ubi perinde est circa quemnam axem fiat incurvatio.
Hic quidem alios casus non sumus contemplati, nisi in quibus ambo
virgae termini vel sunt liberi, vel simpliciter fixi, vel firmiter infixi.
Fieri autem potest ut eadem virga insuper in uno vel pluribus locis
mediis simpliciter figatur, quandoquidem hoc pacto communicatio
inter motus diversarum partium non tolleretur; sed omnium
hujusmodi casuum evolutio requireret tractationem infinitam.
Ut autem ratio calculum ad hujusmodi casum applicandi, intelligatur,
sufficiet unum talem casum hic subjunxisse

\subsection{Problema}
\note{« (150) » en haut, au milieu. Entre « duplex » et « potissimum », un trait vertical à la plume, au début de la ligne. Les lettres \emph{c} et \emph{b} sont soulignées. Le « et » de « et si virga incurvetur » est repris à la plume. Après « fuerit », une tache d'encre couvre une ou deux lettres, et le mot qui suit, lu « parest » avec doute, ne fait pas sens ; il est suivi d'un point, et « Circulus » commence la ligne suivante par une majuscule. « Problema » est souligné, en bas du feuillet, sans texte à la suite.}

\page{193}
(151)

\marginal{fig. 7.}
§. 47. Si virga elastica non solum in utroque termino $E$ et $F$ fuerit simpliciter
fixa, sed etiam in puncto quocunque medio $L$ ope styli figatur, investigare
omnes motus quibus ista virga contremiscere potest.

\emph{Solutio}

Maneat virgae tota longitudo $EF=a$, ac vocetur portio $EL=\lambda a$, ita
ut $\lambda$ denotare possit fractionem quamcunque unitate minorem. Ac primo
quidem patet, si virga in puncto $L$ firmiter esset infixa, omnem
plane communicationem inter ambas portiones $EL$ et $FL$ tolli,
ita ut utriusque motus a motu alterius neutiquam perturbetur.
Verum si in puncto $L$ tantum stylo figatur, circa quem virga
gyrari possit, tum neutra pars motum recipere potest, quin cum altera
is certo modo communicetur. Interim tamen hoc stilo continuitas curvae
per ambas portiones interrumpitur, ita, ut portio $EmL$ alia aequatione
exprimatur atque altera portio $LnF$: scilicet dum hic etiam principio
tantum motus regulares investigamus, qui conformes sunt pendulo simplici $=k$,
pro motu utriusque portionis primus factor $C\sin(\zeta+t\frac{\surd g}{k})$ necessario idem manere
debet, quoniam ambae portiones suas vibrationes eodem tempore similique
modo peragere debent. Alteri autem factores diversi esse \uncertain{poterunt} ratione
coefficientium $\alpha, \beta, \gamma, \delta$. Omisso igitur primo factore ut hactenus faciamus
$\frac{a}{f}=\omega$ et $\frac{s}{a}=u$; tum vero pro portione $EL$ statuamus
\[ y=\ldots\ \alpha e^{u\omega}+\beta e^{-u\omega}+\gamma\sin u\omega+\delta\cos.u\omega \]
pro altera portione $LF$ statuamus
\[ y=\ldots\ \alpha' e^{u\omega}+\beta' e^{-u\omega}+\gamma'\sin u\omega+\delta'\cos.u\omega \]
qui coefficientes a prioribus utcunque discrepare possunt, dummodo observetur,
\note{« (151) » en haut, au milieu ; le nombre 95 est à l'encre en haut à droite. En marge de la première ligne, « fig. 7. », le « 7 » repris à la plume sur un autre chiffre. Une tache d'encre sur « stylo », dont la fin est noyée. « possit » est repris à la plume. Le mot lu « poterunt » est écrit par-dessus une première forme qui se termine en « -erat » ou « -erant ». Dans le second $y$, « cos » est empâté par une surcharge. Un trait vertical après « potest. », à la fin de l'énoncé. « Solutio » est souligné. Le feuillet s'arrête sur « observetur, ».}

\page{194}
(152)

pro puncto $L$, ubi fit $u=\lambda$, ex utraque formula eosdem valores tam
pro $\left(\frac{dy}{ds}\right)$ quam pro $\left(\frac{ddy}{ds^2}\right)$ prodire debere, quandoquidem anguli, quos
utraque portio in $L$ cum axe facit, necessario aequales esse debent; neque
vero etiam radius osculi in hoc puncto $L$ utrinque diversus esse potest.
His observatis pro hoc puncto $L$, ubi fit $u=\lambda$, quia utraque applicata
$y$ evanescere debet, habebimus sequentes quatuor aequationes:
\begin{align*}
&\text{I.}\ \alpha e^{\lambda\omega}+\beta e^{-\lambda\omega}+\gamma\,\text{sin.}\,\lambda\omega+\delta\cos.\lambda\omega=0 \\
&\text{II.}\ \alpha' e^{\lambda\omega}+\beta' e^{-\lambda\omega}+\gamma'\,\text{sin.}\,\lambda\omega+\delta'\cos.\lambda\omega=0 \\
&\text{III.}\ \alpha e^{\lambda\omega}-\beta e^{-\lambda\omega}+\gamma\cos.\lambda\omega-\delta\,\text{sin.}\,\lambda\omega=\alpha' e^{\lambda\omega}-\beta' e^{-\lambda\omega}+\gamma'\cos\lambda\omega-\delta'\,\text{sin.}\,\lambda\omega. \\
&\text{III.}\ (\alpha-\alpha')e^{\lambda\omega}-(\beta-\beta')e^{-\lambda\omega}+(\gamma-\gamma')\cos.\lambda\omega-(\delta-\delta')\,\text{sin.}\,\lambda\omega=0 \\
&\text{IV.}\ (\alpha-\alpha')e^{\lambda\omega}+(\beta-\beta')e^{-\lambda\omega}-(\gamma-\gamma')\,\text{sin.}\,\lambda\omega-(\delta-\delta')\cos.\lambda\omega=0
\end{align*}
Nunc igitur ad utrumque quoque terminum spectamus, unde etiam
quatuor resultabunt aequationes, prouti fuerit vel $\lambda=0$ vel $\lambda=1$;
terminus scilicet $E$, ubi $u=0$, praebet:
\[ \text{V.}\ \alpha+\beta+\delta=0 \quad\text{et}\quad \text{VI.}\ \alpha+\beta-\delta=0, \]
terminus autem $F$, ubi $u=1$, dat.
\[ \text{VII.}\ \alpha' e^{\omega}+\beta' e^{-\omega}+\gamma'\,\text{sin.}\,\omega+\delta'\cos.\omega=0 \quad\text{et} \]
\[ \text{VIII.}\ \alpha' e^{\omega}+\beta' e^{-\omega}-\gamma'\,\text{sin.}\,\omega-\delta'\cos.\omega=0. \]
Nacti scilicet sumus octo aequationes pro definiendis octo coefficientibus,
$\alpha, \beta, \gamma, \delta$ et $\alpha', \beta', \gamma', \delta'$; ac tum supererit adhuc aequatio, unde
angulum $\omega$ definiri oportebit. Incipiamus ab aequatione V et VI,
ex quibus statim colligitur $\beta=-\alpha$ et $\delta=0$; deinde VII et VIII, invicem
additae dant $\alpha' e^{\omega}+\beta' e^{-\omega}=0$ subtractae vero dant
\note{« (152) » en haut, au milieu. Entre « utraque » et « portio », un trait vertical à la plume. Une tache sur « puncto » dans « pro hoc puncto ». Deux équations portent le numéro « III » : la première égale les deux dérivées, la seconde en fait passer les termes du même côté ; les numéros sont transcrits tels qu'ils sont écrits, et la seconde équation est laissée sans être confrontée à la première. Dans « prouti fuerit vel $\lambda=0$ », le premier « vel » est empâté par une surcharge ; les valeurs $0$ et $1$ sont données à $\lambda$ sur le feuillet, bien que la suite parle de $u$. « additae » est écrit sur une première forme, empâtée. Le feuillet s'arrête sur « dant ». Cette page est photographiée une seconde fois à la vue 196.}

\page{195}
(153)

\[ \gamma'\sin\omega+\delta'\cos\omega=0, \text{ unde fit} \]
\[ \beta'=-\alpha' e^{2\omega} \text{ et } \delta'=-\gamma'\,\text{tang.}\,\omega. \]
Hi valores in primo et secundo substituti praebent
\[ \text{I.}\ \alpha(e^{\lambda\omega}-e^{-\lambda\omega})+\gamma\,\text{sin.}\,\lambda\omega=0 \text{ et} \]
\[ \text{II.}\ \alpha'(e^{\lambda\omega}-e^{\omega(2-\lambda)})+\gamma'\,\text{sin.}\,\lambda\omega-\gamma'\cos.\lambda\omega\,\text{tang.}\,\omega=0, \]
ex quibus aequationibus reperiuntur valores:
\[ \gamma=-\frac{\alpha(e^{\lambda\omega}-e^{-\lambda\omega})}{\text{sin.}\,\lambda\omega} \text{ et } \gamma'=-\frac{\alpha'(e^{\lambda\omega}-e^{\omega(2-\lambda)})}{\text{sin.}\,\lambda\omega-\cos.\lambda\omega\,\text{tang.}\,\omega} \]
unde fit
\[ \delta'=\frac{\alpha'(e^{\lambda\omega}-e^{\omega(2-\lambda)})\,\text{tang.}\,\omega}{\sin\lambda\omega-\cos\lambda\omega\,\text{tang.}\,\omega}. \]
Ex his valoribus nunc pro reliquis aequationibus colligimus: $\beta-\beta'=-\alpha+\alpha' e^{2\omega}$
et
\[ \gamma-\gamma'=-\frac{\alpha(e^{\lambda\omega}-e^{-\lambda\omega})}{\text{sin.}\,\lambda\omega}+\frac{\alpha'(e^{\lambda\omega}-e^{\omega(2-\lambda)})}{\text{sin.}\,\lambda\omega-\cos.\lambda\omega\,\text{tang.}\,\omega} \text{ et} \]
\[ \delta-\delta'=-\frac{\alpha'(e^{\lambda\omega}-e^{\omega(2-\lambda)})\,\text{tang.}\,\omega}{\text{sin.}\,\lambda\omega-\cos.\lambda\omega\,\text{tang.}\,\omega} \]
Nunc autem fit
\begin{align*}
&\text{III.}\ \alpha\big((e^{\lambda\omega}+e^{-\lambda\omega})-(e^{\lambda\omega}-e^{-\lambda\omega})\cot.\lambda\omega\big) \\
&\qquad=\alpha'\Big((e^{\lambda\omega}+e^{\omega(2-\lambda)})-(e^{\lambda\omega}-e^{\omega(2-\lambda)})\frac{(\cos.\lambda\omega+\text{sin.}\,\lambda\omega\,\text{tang.}\,\omega)}{\text{sin.}\,\lambda\omega-\cos.\lambda\omega\,\text{tang.}\,\omega}\Big)
\end{align*}
Quae aequatio contrahitur in sequentem formam:
\begin{align*}
&\alpha\big((e^{\lambda\omega}+e^{-\lambda\omega})-(e^{\lambda\omega}-e^{-\lambda\omega})\cot.\lambda\omega\big) \\
&\qquad=\alpha'\big((e^{\lambda\omega}+e^{\omega(2-\lambda)})+(e^{\lambda\omega}-e^{\omega(2-\lambda)})\cdot\cot(1-\lambda)\omega\big).
\end{align*}
quarta vero aequatio praebet istam:
\[ 2\alpha(e^{\lambda\omega}-e^{-\lambda\omega})=2\alpha'(e^{\lambda\omega}-e^{\omega(2-\lambda)}); \]
\note{« (153) » en haut, au milieu ; le nombre 96 est à l'encre en haut à droite. Deux petits traits en marge gauche, en regard de « Hi valores » et de « et $\gamma-\gamma'$ », sans texte. Dans l'égalité III, la parenthèse fermante du membre de droite est écrite à l'intérieur du numérateur de la fraction ; elle est rendue à la fin de l'expression. Le bas du feuillet est blanc sous la dernière équation. Cette page est photographiée une seconde fois à la vue 197.}

\page{198}
(154)

ex hac postrema deducitur
\[ \frac{\alpha}{\alpha'}=\frac{e^{\lambda\omega}-e^{\omega(2-\lambda)}}{e^{\lambda\omega}-e^{-\lambda\omega}}, \]
unde capere licebit
\[ \alpha=e^{\lambda\omega}-e^{\omega(2-\lambda)} \text{ et } \alpha'=e^{\lambda\omega}-e^{-\lambda\omega}, \]
quos valores jam in tertia aequatione substitui oportet: unde facta
reductione prodibit
\[ 0=2-2e^{2\omega}-(e^{\lambda\omega}-e^{\omega(2-\lambda)})(e^{\lambda\omega}-e^{-\lambda\omega})(\cot.\lambda\omega+\cot(1-\lambda)\omega) \]
sive
\[ 0=2-2e^{2\omega}-(e^{\lambda\omega}-e^{\omega(2-\lambda)})(e^{\lambda\omega}-e^{-\lambda\omega})\frac{\text{sin.}\,\omega}{\sin\lambda\omega\,\text{sin.}(1-\lambda)\omega}. \]
Sicque totum negotium huc est perductum, ut ex ista aequatione
valores anguli $\omega$ eliciantur, quem quidem laborem, ob formulas
tantopere complicatas, suscipere non \uncertain{ausim}, si quidem hoc loco
sufficiat methodum tradidisse, qua hujusmodi quaestiones arduae
sint tractandae. Caeterum patet; si fuerit $\lambda=0$, ut portio
$EL$ fiat infinite parva, nullum motum communem inter ambas
partes existere posse, id quod etiam calculus ostendit, quippe
qui praebet
\[ 0=2(1-e^{2\omega}), \]
unde sequitur $e^{2\omega}=1$, ideoque $\omega=0$, quo valore status
quietis innuitur, quod idem eveniet, si statuatur $\lambda=1$, tum
enim terminus $F$ erit quasi muro infixus. At casus, quo
$\lambda=\frac{1}{2}$, singularem evolutionem meretur.
\note{« (154) » en haut, au milieu. Le mot lu « ausim » a une lettre médiane haute qui ressemble à un « t ». Sous « aequatione », « sequitur » et « ideoque », de petits chevrons tracés à la plume, sans addition au-dessus. Le bas du feuillet est blanc.}

\page{199}
(155)

\section{Evolutio casus, quo virga elastica EF non solum in utroque termino E et F, sed etiam in ejus medio L stylo est fixa}

\marginal{Tab. III fig. 8}
§. 48. Cum igitur hoc casu sit $\lambda=\frac{1}{2}$, aequatio finalis hanc induet
formam:
\[ 0=2-2e^{2\omega}-(e^{\frac{1}{2}\omega}-e^{\frac{3}{2}\omega})(e^{\frac{1}{2}\omega}-e^{-\frac{1}{2}\omega})\frac{\sin\omega}{\text{sin.}^2\frac{1}{2}\omega}, \]
quae reducitur ad hanc:
\[ 0=2(1-e^{2\omega})\,\text{sin.}^2\tfrac{1}{2}\omega-(1-e^{\omega})(e^{\omega}-1)\,\text{sin.}\,\omega, \]
quae per factorem communem $1-e^{\omega}$ divisa, (quippe ex quo oriuntur
$e^{\omega}=1$, hinc $\omega=0$) pro statu quietis producit hanc aequationem:
\[ 0=2(1+e^{\uncertain{2\omega}})\sin^2\tfrac{1}{2}\omega-(e^{\omega}-1)\,\text{sin.}\,\omega, \]
quae, si loco $\text{sin.}\,\omega$ scribamus $2\,\text{sin.}\,\frac{1}{2}\omega\cos.\frac{1}{2}\omega$, abit in hanc:
\[ 0=(1+e^{\omega})\sin^2\tfrac{1}{2}\omega-(e^{\omega}-1)\sin\tfrac{1}{2}\omega\cos.\tfrac{1}{2}\omega, \]
quae manifesto duos habet factores, alterum $\text{sin.}\,\frac{1}{2}\omega$, alterum vero
\[ (1+e^{\omega})\sin\tfrac{1}{2}\omega-(e^{\omega}-1)\cos\tfrac{1}{2}\omega, \]
quorum uterque nihilo aequatus praebet solutionem: ambas
igitur seorsim perpendamus.

§. 49. Pro priore igitur casu statuamus $\sin\frac{1}{2}\omega=0$, eritque
in genere $\frac{1}{2}\omega=i\pi$, denotante $i$ numerum integrum quemcunque; ita
ut hic jam innumerabiles motus regulares contineantur; ac manifestum
est hunc casum prorsus convenire cum casu quarto supra evoluto, nisi
\note{« (155) » en haut, au milieu ; le nombre 97 est à l'encre en haut à droite. Le titre « Evolutio casus, » est souligné, et la dernière ligne du titre d'un long trait ; il ne porte pas de numéro. L'exposant du premier $e$ dans la troisième équation est surchargé et empâté : « $2\omega$ » est la lecture la plus probable de ce qui est visible. Un petit signe isolé, en forme de « o », dans la marge devant la deuxième équation. Une tache d'encre sous « denotante ». Le feuillet s'arrête sur « nisi ».}

\page{200}
(156)

quod hic sit $\frac{1}{2}\omega$ quod ibi erat $\omega$; scilicet hic pro utraque
portione longitudo tantum est $\frac{1}{2}a$. Utraque igitur semissis eodem
modo suas vibrationes peragit, ac si seorsim existeret et in utroque
termino stylo simpliciter esset fixa; quamobrem omnes soni simplices,
quos utraque portio edere potest, sequentibus numeris exprimentur
\[ \frac{4\pi c\surd 2gb}{aa};\ \frac{16\pi c\surd 2gb}{aa};\ \frac{36\pi c\surd 2gb}{aa};\ \frac{64\pi c\surd 2gb}{aa}; \]
qui ergo omnes duplici octava altiores sunt quam casu IV;
cujus discriminis ratio in hoc est sita, quod longitudo hoc
casu tantum semissis est illius.

§. 50. Hoc igitur casu coefficientes $\alpha, \beta, \gamma, \delta$. et $\alpha', \beta', \gamma', \delta'$,
sequenti modo determinabuntur:
\[ \alpha=e^{\frac{1}{2}\omega}-e^{\frac{3}{2}\omega};\quad \beta=e^{\frac{3}{2}\omega}-e^{\frac{1}{2}\omega};\quad \gamma=-\frac{(e^{\frac{1}{2}\omega}-e^{\frac{3}{2}\omega})(e^{\frac{1}{2}\omega}-e^{-\frac{1}{2}\omega})}{\sin\frac{1}{2}\omega};\quad \delta=0 \]
\[ \alpha'=e^{\frac{1}{2}\omega}-e^{-\frac{1}{2}\omega};\quad \beta'=-e^{2\omega}(e^{\frac{1}{2}\omega}-e^{\frac{1}{2}\omega});\quad \gamma'=-\frac{(e^{\frac{1}{2}\omega}-e^{-\frac{1}{2}\omega})(e^{\frac{1}{2}\omega}-e^{\frac{3}{2}\omega})}{\sin\frac{1}{2}\omega-\cos\frac{1}{2}\omega\,\text{tang.}\,\omega} \text{ et} \]
\[ \delta'=\frac{(e^{\frac{1}{2}\omega}-e^{-\frac{1}{2}\omega})(e^{\frac{1}{2}\omega}-e^{\frac{3}{2}\omega})\,\text{tang.}\,\omega}{\sin\frac{1}{2}\omega-\cos\frac{1}{2}\omega\,\text{tang.}\,\omega}. \]
Est vero
\[ \sin\tfrac{1}{2}\omega-\cos\tfrac{1}{2}\omega\,\text{tang.}\,\omega=\frac{\sin\frac{1}{2}\omega\cos.\omega-\cos.\frac{1}{2}\omega\,\text{sin.}\,\omega}{\cos.\omega}=-\frac{\sin\frac{1}{2}\omega}{\cos.\omega} \]
ubi, quia est $\omega=2i\pi$, erit $\cos.\omega=1$, at \struck{\ill{}} $\text{sin.}\,\frac{1}{2}\omega=0$
Multiplicemus igitur omnes hos coefficientes per $\sin\frac{1}{2}\omega$ eritque
\[ \alpha=0,\quad \beta=0,\quad \gamma=-(1-e^{\omega})(e^{\omega}-1) \quad\text{et}\quad \delta=0; \]
\note{« (156) » en haut, au milieu. Le chiffre du numérateur des exposants $\frac{3}{2}$ a une boucle fermée qui le rapproche d'un « 9 », comme le « 3 » de la vue 186 ; il est lu « 3 », valeur qu'il a à la vue 199. Dans $\beta'$, les deux exposants de la parenthèse sont écrits $\frac{1}{2}\omega$ tous les deux, et ils sont laissés tels. Les nombres des sons, « 36 » et « 64 », sont écrits d'un « 6 » à boucle ouverte qui ressemble à un « b ». Devant « sin. $\frac{1}{2}\omega = 0$ », un mot court est biffé d'une rature épaisse et n'est pas lu. Le « l » de « longitudo », à la ligne « quod longitudo hoc », est repris à la plume. Le feuillet s'arrête sur « $\delta=0$; » ; la suite est à la vue 201.}

\end{document}
